\chapter{Conclusion \& Perspective}

This thesis documented the design and implementation of Anhoc, a web-based mathematics learning platform that combines bilingual lesson content, generated practice exercises, grade-level assessments, progression tracking, and administrative control. The implemented system uses a modern web stack with a Next.js frontend, an Express backend, Prisma as the persistence layer, and PostgreSQL as the database.

The project demonstrates several sound software engineering decisions. Most importantly, it separates reusable question templates from persisted question snapshots, enabling reproducible grading and better issue reporting. It also uses a dynamic access-control model and a modular progression pipeline for mastery, XP, and achievements. These choices make the system more extensible than a minimal educational CRUD application.

At the same time, the current implementation remains an evolving platform rather than a finished production learning system. It still needs broader automated testing, stronger validation of author-authored template logic, clearer schema ownership, and empirical studies of educational effectiveness. Even with these limitations, the project is a solid foundation for future academic work and practical product development in mathematics learning software.
